\section{The supplied Bessel wave and the choice of spectral data}

Fix an original datum $D$, a setting $X$ over $D$, and, where it is
needed, $Y:\mathrm{Full}(X)$. Two spaces occur. The space
$\mathcal H_X=\mathrm{HHX}(X)$ is the image of the finite and infinite
operators on the adelic automorphic quotient. The real Bessel quotient
for a word $u$ is $S_{u,\infty}\backslash G_\infty$, with sections
transforming on the left by $\psi_{u,\infty}\pi_u$. The Bessel wave
is a section on this second quotient. We will use a dimension bound
on each of the two spaces, for different purposes.

% BEGIN SFS1
\paragraph{SFS1. Finite dimension of the actual automorphic range.}
\label{claim:SFS1}
If the original infinite kernel $X.\mathrm{ph}$ is integrable,
$\dim_{\mathbb C}\mathcal H_X$ is finite. There is no gap condition
in this statement.

\begin{proof}
At each real place let $w_v$ be the given lowest unit vector,
$\sigma_v$ the given discrete series, $\mu_v$ the prescribed
group measure, and put
\[
 d_v=\left(\int|\langle w_v,\sigma_v(g)w_v\rangle|^2\,d\mu_v(g)
       \right)^{-1},\qquad
 p_v(g)=d_v\overline{\langle w_v,\sigma_v(g)w_v\rangle}.
\]
This is $X.\mathrm{dv}(v)$: the numerator in the source's
formal-degree formula is $\|w_v\|^4=1$. Discrete-series Schur
orthogonality gives a positive formal degree for $\mu_v$.
Specializing all four vectors in that identity to $w_v$ shows
that this degree equals $d_v$. Thus its four-vector identity holds
with precisely the scalar above. In particular $p_v$ is continuous
and $p_v(1)=d_v>0$.

By definition $X.\mathrm{ph}(g)=\prod_v p_v(g_v)$ and the
infinite measure is the product of the $\mu_v$. Fix one place $v$.
At every other place choose a relatively compact identity
neighborhood on which $|p_j|\ge d_j/2$.
Its Haar measure is finite and positive. Tonelli on the product
of those neighborhoods with the entire $v$-factor bounds a
positive constant times $\int |p_v|$ by $\int |X.\mathrm{ph}|$.
For a single place the constant is one. Hence every local kernel
has the $L^1$ property required in RFD8.

Apply RFD8 to $X.B$'s adelic domain, measure, right regular action,
finite test function, and finite convolution operator. Local
constancy and support in the lattice stabilizer are fields of its
admissible finite test. The local representations are irreducible
and unitary by the discrete-series field. The previous paragraph
supplies the kernel hypotheses. The weak infinite convolution
operator is exactly $X.R_{\mathrm{ph}}$ by its defining field.
The range in RFD8 is therefore $\mathcal H_X$ on the supplied quotient.
\end{proof}
% END SFS1

% BEGIN SFS2
\paragraph{SFS2. Positivity of both original word volumes.}
\label{claim:SFS2}
For every supplied $X,Y,u$ the level normalization holds and
\[
 0<Y.\mathrm{volN}(u)=\mu_f(K_2)_{\mathbb R}\,
                  \mu_{N,u}(Y.FN(u))_{\mathbb R},\qquad
 0<Y.\mathrm{volH}(u)=\mu_{H,u}(G_\infty)_{\mathbb R}.
\]
Here the real conversion of an infinite measure is zero by source
convention, so actual finiteness of each factor is part of the proof.

\begin{proof}
ASG14 applies to the reference frame of $X$: its residue
congruences place $K_g$ in $K_2$ with trivial character, and place
$K_2$ inside the lattice stabilizer. The former is open and the
latter compact. Since $K_2$ is a subgroup, it is compact open
and has positive finite mass for the given finite Haar measure.
These are the four original normalization clauses.

For $N=N_{u,\infty}$ use exactly the rational group
$L=\mathrm{NOinf}(X,u)$: at the finite places its elements
belong to $N_u(\mathbb A_f)\cap x_uK_2x_u^{-1}$. This is a
compact open subgroup of $N_u(\mathbb A_f)$. In particular $L$
contains all the integral exponentials needed for the compact
cover in FLD2. In detail, for a rational Malcev basis $Z_i$ as
in FLD2, clear the denominators of the exponential polynomials
and then multiply each $Z_i$ by a common multiple of the
congruence primes to the powers required by that finite open
subgroup. At the finitely many exceptional places the exponential
of a sufficiently divisible integral multiple approaches one
$p$-adically. At all other places integrality suffices. Both
$\exp(\pm Z_i)$ after these scalings belong to $L$. The decreasing
coordinate reduction of FLD2 gives a compact $C\subset N$
with $LC=N$. Also $L$ is countable, being an image of rational points.

The measure $\mu_{N,u}$ is supported on $N$ and restricts there
to a nonzero Haar measure finite on compact sets by the
sigma-finite invariance argument of FLD2. Put $A=Y.FN(u)\cap N$.
The supplied strict fundamental-domain field, restricted to $N$,
makes $A$ measurable with one representative per $L$-orbit.
If it were null then so would be $N=\bigcup_{l\in L}lA$.
To bound its measure above, enumerate $L$ and partition $A$ by
the first $l$ for which $la\in C$. The translated pieces are
disjoint by strictness, so $\mu_{N,u}(A)\le\mu_{N,u}(C)<\infty$.
This proves the assertion on the very domain stored in $Y$.
At infinity the residual $H_{u,\infty}$ is compact by PGR2 and
the definite orthogonal complement of PGR5. Its given measure
restricts to nonzero Haar measure by FLD2, hence its total mass
is finite and positive. Taking real values proves the formulas.
\end{proof}
% END SFS2

% BEGIN SFS3
\paragraph{SFS3. Multiplicity of the local projector on Bessel sections.}
\label{claim:SFS3}
Use any one of the actual real places and words under the original
$\mathrm{GapHyp}(D)$. Write $G_v$ for its real unitary group,
$S_v=N_v\rtimes H_v$ for its Bessel subgroup,
$\rho_v=\psi_v\pi_v$ for the given unitary action on $V_v$,
and $\mathcal I_v$ for the Hilbert space of sections with
$f(sx)=\rho_v(s)f(x)$ and finite squared norm on $S_v\backslash G_v$.
Convolution in the right translation action $R_v$ by $p_v$ of
SFS1 is an orthogonal projection of rank at most one on
$\mathcal I_v$.

\begin{proof}
BWP6, BWP20 and BWP23 give the integrability of $p_v$ for the
supplied Haar measure and lowest vector. RFD5 applies to every
strongly continuous unitary action. Here the quotient is unimodular,
and translations act unitarily; continuity follows by density of
continuous compactly supported sections. Write $P_v=R_v(p_v)$.
For $b\in\operatorname{ran}P_v$, RFD6 constructs the bounded
intertwiner $A_b$ from the Hilbert representation $\sigma_v$ into
$\mathcal I_v$ by $A_b(\sigma_v(g)w_v)=R_v(g)b$. The Schur
compression formula gives $\|A_b\|=\|b\|$. Thus evaluation at
$w_v$ in this construction distinguishes the intertwiners attached
to distinct vectors in the image.

Every bounded such intertwiner $A$ defines a continuous map on
the space of smooth vectors of the discrete series by
\[
                     B_A(a)=(Aa)(1)\in V_v.
\]
Here evaluation requires justification, because $Aa$ initially
denotes an $L^2$ class. A smooth vector for the action on
$\mathcal I_v$ has weak derivatives of every order in the
infinitesimal translation directions. In a relatively compact
local trivialization of the smooth homogeneous bundle, these
directions span the quotient tangent space by transitivity.
Solve for coordinate derivatives in terms of a spanning subfamily
on a smaller chart. Bundle transition and vertical action
coefficients are smooth, with bounded derivatives on compact
subcharts. Induction, with a compact cutoff, bounds the local
Sobolev norm of order $j$ by a constant times a finite sum of the
Hilbert norms of translation derivatives of orders at most $j$.
The invariant quotient measure has a smooth positive density in
this chart. Local Sobolev embedding for $j>\dim(S_v\backslash
G_v)/2$ gives a continuous evaluation; taking all $j$ gives a
unique smooth representative. Bounded intertwining commutes
with all these Hilbert derivatives, so $B_A$ is continuous for
the smooth-vector topology of $\sigma_v$. It obeys
$B_A(\sigma_v(s)a)=\rho_v(s)B_A(a)$. If $B_A=0$, evaluating
this construction on $\sigma_v(g)a$ gives $(Aa)(g)=0$ for
every $g$ and smooth $a$. Density then gives $A=0$.

For the resulting multiplicity bound we use the published
archimedean Bessel uniqueness theorem of D. Jiang, B. Sun and
C.-B. Zhu, \emph{Uniqueness of Bessel models: the archimedean
case}, Theorem A and its flag description in Section 2
(\href{https://arxiv.org/abs/0908.1728}{arXiv:0908.1728}).
In its notation the successive isotropic spaces are the spans
of the first $j$ vectors of the given $e$-frame; adjoining the
negative line generated by $z_u$ gives the last flag space.
The ambient signature is $(k+r,r+1)$, and the residual group is
the compact $U(k)$ on the positive complement. Each simple
additive quotient $L_j$ is a complex line. BWP20 computes on it
the additive character with nonzero coefficient $4\pi/g_D$,
including the last link to $z_u$. Its restriction to every
$L_j$ is therefore nontrivial, as genericity requires. It is unitary, and is trivial on the
compact part in the scalar-character formulation.
The representation on the large group is the smooth
globalization of the given irreducible discrete series, and the
second representation in the theorem is the contragredient
of the given irreducible compact type $V_v$. These are irreducible
Casselman--Wallach representations. The functionals to which the
theorem applies are
$(a,\ell)\mapsto\ell(B_A(a))$ on
$\sigma_v^\infty\widehat\otimes V_v^*$.
They transform by exactly $\psi_v$ since the two compact actions
cancel. The finite-dimensional dual separates the values of
$B_A$, and evaluation distinguished the intertwiners above.
Theorem A bounds this Hom space by dimension one. This proves
the assertion for the original coefficient projector. The
bound concerns its image on Bessel sections for this specified
representation and word.
\end{proof}
% END SFS3

% BEGIN SFS4
\paragraph{SFS4. Decay and compact reproduction for the supplied wave.}
\label{claim:SFS4}
The full original $\mathrm{WuDecayOfGap}$ holds at the published
multiplicity-one input of SFS3. For every $D$ satisfying the
original gap and \emph{every} supplied $X$ over it, the actual
functions $X.W_u$ satisfy weighted integrability and compact
reproduction for all words, and $X.\mathrm{ph}$ is integrable.

\begin{proof}
Fix the word. The real stabilizer, character and compact
representation on the full tensor space are respectively the
products of the local stabilizers, characters and compact
actions. For the representations this follows from the supplied
tensor inner-product, spanning and compact equivariance fields,
and the trivial action on the unipotent radical. PGR2 gives
the semidirect decomposition. BWP19 identifies the pushforward
of the given quotient measure to $S_\infty\backslash G_\infty$
with a finite positive multiple of the product of local quotient
Haar measures. Its proof uses a cutoff and uniqueness of the
quotient measure, so the identity extends from product test
functions to the whole Borel sigma-algebra. In particular the
Hilbert space of global Bessel sections is, up to an overall
norm factor, $\widehat\bigotimes_v\mathcal I_v$.
One may see surjectivity in measurable local trivializations:
product simple functions of finite support in measure and pure
fiber vectors span a dense subspace. This identification intertwines
the componentwise translations; the common scalar in the norm
does not change their integrated action.

By SFS3, $p_v\in L^1$ at every place. Finite-product Fubini gives
$X.\mathrm{ph}=\prod_v p_v\in L^1$ and the integrated
operator on sections equals $P=\widehat\bigotimes_v P_v$.
Indeed its weak matrix coefficients agree on pure tensors,
hence everywhere by the $L^1$ bound. Therefore
$\dim\operatorname{ran}P\le1$. The supplied function $W=X.W_u$
is already a normalized continuous equivariant section by
$X.Wu\_spec$. Its pointwise convolution by $X.\mathrm{ph}$
equals $W$. This equality also holds in the Hilbert space:
pair with an $L^2$ section $b$ and apply Fubini, using
$\int |X.\mathrm{ph}(g)|\,\|b\|_2\|R(g)W\|_2\,dg<\infty$.
Thus $PW=W$ and $\operatorname{ran}P$ is a line.

Use separately the packet $W^0$ of BWP20--BWP22, with nonvanishing
from BQS23 and the numerical implication BWP23. Their construction
uses the frames, actual compact representation, lowest vectors,
tensor map, and subgroup and quotient measures supplied in this
same $X$. In particular it constructs, for this comparison, a
normalized section for the identical $\rho_u$ and $P$.
Its defining convolution shows $PW^0=W^0$. Hence $W=cW^0$
in the quotient Hilbert space with $|c|=1$ by the common
normalization. Both sections are continuous, and their difference
has a continuous invariant norm. Haar quotient measure has full
support, so the equality holds for every point of the group.

The compactly supported bounded reproducing kernel constructed for
$W^0$ in BWP22 consequently reproduces $W$ pointwise, for the
same infinite Haar measure. Their absolute norms coincide, so
the weighted $L^1$ bound of BWP22 transfers with the literal
$\mathrm{nuInf}^{N_X}$ and $X.\nu_Q(u)$. The scalar may depend
on the word. This proves both clauses word by word together
with integrability of the common infinite kernel.
\end{proof}
% END SFS4

% BEGIN SFS5
\paragraph{SFS5. The full finiteness and compatibility record on $Y$.}
\label{claim:SFS5}
For any $D,X,Y$, the two hypotheses
$\mathrm{WuDecay}(X)$ and $X.\mathrm{ph}\in L^1$ imply
$\mathrm{FinX}(Y)$. In particular the record exists on every
supplied tied pair $\mathrm{ConcTie}(X,Y)$ without an arithmetic
gap hypothesis. Under $\mathrm{GapHyp}(D)$ it exists for every
$X,Y$. This gives the original $\mathrm{FinXC}$ for every
spectral tie over such a $D$.

\begin{proof}
For each word PGR6 identifies the original map
$\mathrm{sharp}$ with evaluation on the supplied finite column
representatives followed by Riesz duality. Its explicit inverse
assigns to a representative the corresponding functional and
extends by covariance; PGR5 makes this extension unique.
This supplies both injectivity and surjectivity, with their actual
complex-antilinear convention. SFS2 and FLD3 give folding on
every word. QPR3 uses $X.\mathrm{ph}\in L^1$ and those folds to
put the Poincar\'e functions in $\mathcal H_X$. For the particular
transversal $Y.Dt(u)$, the seed for a sum of two columns is the
sum of their seeds by PGR7 and the defining linear evaluation;
scalar multiplication has the same compatibility. Absolute
summability pointwise from the fold/range statement allows
addition and scalar multiplication through the period series.
Passing to their $L^2$ classes respects these operations. These
are the original $\mathrm{inHH}$ and $\mathrm{puLin}$.

SFS1 gives finite dimension of $\mathcal H_X$. ARF1 supplies
finiteness over each arithmetic term; ARF2--ARF3 give independence
of both the pair corank and its scalar predicate on equal-row
lifts. ARF4 puts every Gram parameter in the real fixed subfield
and supplies the one positive signature scaling for the Haar
comparison. SFS2 gives positive real volumes for both factors
at each word. These are all ten fields of $\mathrm{FinX}(Y)$.
The tied specialization takes the two analytic hypotheses from
the original $\mathrm{ConcTie}$. For the gap specialization
take them from SFS4 on the given $X$, and for
$\mathrm{FinXC}$ use $X=T.X$, $Y=T.Y$ for the given spectral tie.
\end{proof}
% END SFS5

% BEGIN SFS6
\paragraph{SFS6. The local comparisons and the concrete spectral package.}
\label{claim:SFS6}
Suppose $T:\mathrm{STies}(D)$ is supplied. If $T.X$ satisfies
the two analytic hypotheses of SFS5 and
\[
   1\le T_w\left(\frac{D.\mathrm{host}.\mathrm{ka}}
                                {\mathrm{gg}(D)}\right)^{2r}
            \quad\text{for every word }w,\qquad
            T_w=(D.\mathrm{mem}(w)).A.T,\quad r=D.R.rS ,
\tag{SFS.1}
\]
then the original $\mathrm{STiesOK}(T)$ holds and there is a
$\mathrm{ConcK}(D)$ with $X=T.X$, $Y=T.Y$. The original
$\mathrm{GapHyp}(D)$ suffices for the two analytic hypotheses.

\begin{proof}
Write $h_F$ for SFS5 and specialize the original field of $T$ to
this very witness to obtain $I=T.\mathrm{hI}(h_F)$, the
comparison $I:D.s\simeq s_Y$ with $s_Y=\mathrm{sOf}(h_F)$.
For every finite prime $p$ its local field $I.\mathrm{atP}(p)$
is an isomorphism from $(D.s).\mathrm{atP}(p)$ to
$s_Y.\mathrm{atP}(p)$. The latter is \emph{defined} to be
$\mathrm{locC}(D,p)$. This supplies the entire local comparison
$\mathrm{LocCP.LocTied}(D)$ using all primes, with no choices
of new local models.

The isometry $I.\mathrm{eCu}(w)$ identifies the given abstract
column module with $\mathrm{ColT}(h_F.\mathrm{shInj}(w))$.
The original column comparison $D.s\_Cu(w)$ identifies the rank
of the former with $D.\mathrm{duS}(w)$, so the real-valued
$\mathrm{duC}$ equality follows with these exact spaces.
The equality for $D_\ast$ is the mass-normalization field of
$T.\mathrm{autTie}$. SFS2 gives the normalized level and PGR5
gives neatness for all words. Use the original $T.\mathrm{hb}$,
$T.\psi$ and their supplied measure and character comparisons for
the archimedean and finite components; append the local isomorphisms
just obtained. The two analytic hypotheses and (SFS.1), which is
$\mathrm{TFrameAt}(D)$ verbatim, give the remaining clauses of
$\mathrm{ConcTie}(T.X,T.Y)$. Together these statements also
give each conjunct in $\mathrm{STiesOK}(T)$. The required
concrete package is $(T.X,T.Y,h_F,I,\mathrm{ConcTie})$.
Under the original gap SFS4 applies to the specific setting $T.X$.
In particular QPR11 now applies to every pair of its words, with
the given kernel normalization and the same comparison isometries.
\end{proof}
% END SFS6

% BEGIN SFS7
\paragraph{SFS7. Unconditional and hypothesis-dependent family thresholds.}
\label{claim:SFS7}
Let $\mathcal F$ be an original family with an $\mathrm{SSelF}$
witness $(a,T)$. If there is a number $b$ such that, for every
$x\ge b$ with $x>\mathcal F.x_1$ and $x\ge a$, every
$i\in\mathcal F.Ix(x)$ satisfies $\mathrm{GapHyp}(D_{x,i})$
and (SFS.1), then the family has an original $\mathrm{CSelF}$
witness. The same assertion holds with the two analytic
hypotheses for the selected $T_{x,i}.X$ in place of the gaps.
If this uniform-in-$i$ property holds on every family of a
class with an original $\mathrm{SSel}$, these choices give
$\mathrm{CSel}$ on that class.

There is also a conclusion when the numerical data are supplied
only under the original standing hypotheses. For any family class
with a supplied $\mathrm{SSel}$, a closed original
$\mathrm{PX.TFrame}$ event implies the closed $\mathrm{SelOK}$ event
at the parameter values
\[
 (C_g,c_b,C_u,c_T,C_T)
  =(\mathcal F.\mathrm{CYL}+1,\tfrac14,3,\tfrac12,\mathcal F.\mathrm{CTu}),
\]
for every late $x$ and every $i\in\mathcal F.Ix(x)$ in
$\mathrm{OwnEv}$, with the same $C_0$. It therefore implies the closed
$\mathrm{ConcEv}$ event. This latter implication
applies in particular to the selector induced by any supplied
$\mathrm{TSel}$ and its compatible $\mathrm{SelS}$ in
$\mathrm{link\_alt}$.

\begin{proof}
For the first assertion take $a'=\max(a,b)$. Given $x>\mathcal
F.x_1$, $x\ge a'$ and $i$, use the original $T_{x,i}$ at its
threshold $a$. SFS6 supplies its concrete package. This is the
required dependent function of $x$ and $i$, for all indices
simultaneously. If only existence of $b$ was supplied, choose it
once per family. If the original compatible member records come
from $\mathrm{TSel},\mathrm{SelS}$, project the latter to
$\mathrm{STies}$ at the former's threshold, exactly as in the
definition of the selector in $\mathrm{link\_alt}$. This preserves
both the chosen setting and the family point.

For the conditional event, strengthen the closed parameter
thresholds for $C',C_l$ to at least $29$ by taking their maxima
with the constant function $29$ at the middle threshold level.
A family large for the strengthened thresholds is still large
for the old ones, with the same outer accuracy bound. The exact
$\mathrm{OwnEv}$ uses the constants $sC'(\mathcal F)$ and
$sC_l(\mathcal F)$ fixed by this family. Under the event's
admissibility premise the standing inequalities give
$m\ge2$, $f\ge m^4\ge16$, $L>1$ and
$k\ge C'm^5 f L$. Hence $n=2r+k+1\ge1801$.
For adjacent weight indices, the same standing inequalities,
after cancelling their common shift, give
$b_{j+1}-b_j\ge C_l n^{19} f^5 L\ge16n^{19}$.
These are the original gap conditions. At the late threshold for
$\mathrm{PX.TFrame}$ the same admissibility premise, with
its conjuncts in the order used by $\mathrm{Cnt.Hyp0}$, gives
(SFS.1). Apply SFS6 to the selected $T_{x,i}$ for every $i$,
whenever its own threshold and the strict family threshold hold.
This is $\mathrm{SelOK}$. To obtain $\mathrm{ConcEv}$ enlarge
the event threshold to the maximum of its old value,
$a$, and $\mathcal F.x_1+1$. Above it the selected package
exists whenever the original admissibility premise holds.

The last implication is guarded by that premise; $\mathrm{CSelF}$
requires concrete packages at every late point even when the premise
fails. A common-family $\mathrm{SSel}$ and uniform arithmetic
frame inequalities across its points are retained at the scopes
stated here. The distinct-word, diagonal and counting estimates in
the full spectral comparison require further arguments.
\end{proof}
% END SFS7
